% Notes of the Montreal seminar, 1970 (cote 129, typed part): archive cote 129, site page 37 (PDF page 38).
% typescript. Draft compared: none.
% AI-transcribed by Claude Opus 5.5 (claude-opus-5-5, adaptive thinking, effort high), Message Batches API, batch mtl129_full (msgbatch_01AGnDKHwy6ATCM7ehUa6263), 2026-11-13; not yet checked by eye.
\scanpage{129/37}
\mspage{\hand{II.9}}

2.5. \quad Supposons que $S$ soit un $\mathbb{F}_p$-schéma. Pour tout groupe fini
localement libre sur $S$ on définit alors les morphismes de Frobenius
relatif et de Verschiebung (I,5). On a les critères suivants :

2.5.1. \quad $G$ est infinitésimal $<==>$ $\mathbb{F}_{G/S}$ est localement nilpotent.

2.5.2. \quad $G$ est unipotent $<==>$ $\mathbb{V}_G$ est localement nilpotent.

2.5.3. \quad $G$ est bi-infinitésimal $<==>$ $\mathbb{F}_{G/S}$ et $\mathbb{V}_G$ sont localement
nilpotents.

2.5.4. \quad $G$ est étale $<==>$ $\mathbb{F}_{G/S}$ est un isomorphisme.

2.5.5. \quad $G$ est de type multiplicatif $<==>$ $\mathbb{V}_G$ est un isomorphisme.

Il suffit de prouver les assertions 2.5.1. et 2.5.4., les
autres s'en déduisant, d'après 1.4. par dualité de cartier.

\uline{Démonstration de 2.5.1.}

Soit $\varepsilon : S \to G$ la section unité de $G$ ($G$ est considéré
comme schéma en ``ensembles pointés'' grâce à cette section $\varepsilon$ ). Nous al-%
\margin{X} lons caractériser le ``noyau'' de $\mathbb{F}^n_{G/S}$ pour tout S-schéma $G$ muni d'une
section. On vérifie immédiatement que le noyau de $\mathbb{F}^n_{G/S}$ est défini par
l'idéal engendré par $F^n(J) \subset \underline{O}_G$ , image de l'idéal d'augmentation
par le morphisme $x \to x^{p^n}$ ; désignons par $J_n$ cet idéal. On a évidem-%
ment $J_n \subset J^{p^n}$ . Si $J$ est de type fini, on peut supposer qu'il a moins
de $p^k$ générateurs, pour un entier $k$ convenable, et on vérifie immé-%
diatement que l'on a alors
\[ J^{p^{n-k}} \subset J_n , \]
de sorte que les $J_n$ forment un système cofinal aux $J^{p^n}$ .

\margin{37}
